\section{Derive the special second-byte bounds}
The extra bounds prevent three distinct problems: an encoding longer than the
shortest one for a value, a surrogate value, or a value beyond U+10FFFF. They apply
to the next byte after the lead, not to every remaining continuation.

For lead \texttt{e0}, the high scalar bits are zero. A second byte below
\texttt{a0} would leave the completed value below hexadecimal 0800, where a shorter
encoding is required. With second byte \texttt{a0} and final byte \texttt{80}, the
payload bits instead produce exactly 2048, the first value requiring three bytes.

For lead \texttt{ed}, a second byte \texttt{a0} would put even the lowest possible
completion at hexadecimal d800. Restricting it through \texttt{9f} keeps the result
below the surrogate range. The following continuation again uses \texttt{80--bf}.
An accepted \texttt{ed 9f} prefix must therefore allow a final \texttt{bf}; retaining
\texttt{9f} as a permanent upper bound would wrongly reject that valid sequence.

For lead \texttt{f0}, the next byte must begin at \texttt{90}. Its low six bits are
then at least 16, contributing \(16\cdot4096=65536\), the first four-byte scalar.
For lead \texttt{f4}, the second byte may reach only \texttt{8f}. The greatest
completion, \texttt{f4 8f bf bf}, represents U+10FFFF. Increasing that second byte to
\texttt{90} would already make every completion exceed the scalar limit.

These are encoding rules. The scanner accepts well-formed encodings of unassigned
scalar values and noncharacters as well as ordinary text values. It counts an
encoded U+FEFF instead of silently removing it. Applications may impose separate
content policies; those policies are not part of this validator's byte grammar.
